\documentclass[10pt,a4paper]{amsart}
\usepackage{amsmath,amssymb,amsthm,graphicx}
% Retained minimal preamble for the Stacks Project chapter sites-cohomology.tex (Cohomology on Sites).
% Verbatim excerpt of the chapter preamble: Xy-pic setup (native preamble lines 14-18) and the native
% theorem declarations with the mathematical macros (native preamble lines 154-217). Copied mechanically;
% nothing added, removed or redefined. The wrapper class and the amsmath/amssymb/amsthm/graphicx base
% packages are supplied by the assembler, so the class/externaldocument part of the native preamble is omitted.
% For commutative diagrams we use Xy-pic
\usepackage[all]{xy}

% We use 2cell for 2-commutative diagrams.
\xyoption{2cell}
% Theorem environments.
%
\theoremstyle{plain}
\newtheorem{theorem}[subsection]{Theorem}
\newtheorem{proposition}[subsection]{Proposition}
\newtheorem{lemma}[subsection]{Lemma}

\theoremstyle{definition}
\newtheorem{definition}[subsection]{Definition}
\newtheorem{example}[subsection]{Example}
\newtheorem{exercise}[subsection]{Exercise}
\newtheorem{situation}[subsection]{Situation}

\theoremstyle{remark}
\newtheorem{remark}[subsection]{Remark}
\newtheorem{remarks}[subsection]{Remarks}

\numberwithin{equation}{subsection}

% Macros
%
\def\lim{\mathop{\mathrm{lim}}\nolimits}
\def\colim{\mathop{\mathrm{colim}}\nolimits}
\def\Spec{\mathop{\mathrm{Spec}}}
\def\Hom{\mathop{\mathrm{Hom}}\nolimits}
\def\Ext{\mathop{\mathrm{Ext}}\nolimits}
\def\SheafHom{\mathop{\mathcal{H}\!\mathit{om}}\nolimits}
\def\SheafExt{\mathop{\mathcal{E}\!\mathit{xt}}\nolimits}
\def\Sch{\mathit{Sch}}
\def\Mor{\mathop{\mathrm{Mor}}\nolimits}
\def\Ob{\mathop{\mathrm{Ob}}\nolimits}
\def\Sh{\mathop{\mathit{Sh}}\nolimits}
\def\NL{\mathop{N\!L}\nolimits}
\def\CH{\mathop{\mathrm{CH}}\nolimits}
\def\proetale{{pro\text{-}\acute{e}tale}}
\def\etale{{\acute{e}tale}}
\def\QCoh{\mathit{QCoh}}
\def\Ker{\mathop{\mathrm{Ker}}}
\def\Im{\mathop{\mathrm{Im}}}
\def\Coker{\mathop{\mathrm{Coker}}}
\def\Coim{\mathop{\mathrm{Coim}}}

% Boxtimes
%
\DeclareMathSymbol{\boxtimes}{\mathbin}{AMSa}{"02}

%
% Macros for moduli stacks/spaces
%
\def\QCohstack{\mathcal{QC}\!\mathit{oh}}
\def\Cohstack{\mathcal{C}\!\mathit{oh}}
\def\Spacesstack{\mathcal{S}\!\mathit{paces}}
\def\Quotfunctor{\mathrm{Quot}}
\def\Hilbfunctor{\mathrm{Hilb}}
\def\Curvesstack{\mathcal{C}\!\mathit{urves}}
\def\Polarizedstack{\mathcal{P}\!\mathit{olarized}}
\def\Complexesstack{\mathcal{C}\!\mathit{omplexes}}
% \Pic is the operator that assigns to X its picard group, usage \Pic(X)
% \Picardstack_{X/B} denotes the Picard stack of X over B
% \Picardfunctor_{X/B} denotes the Picard functor of X over B
\def\Pic{\mathop{\mathrm{Pic}}\nolimits}
\def\Picardstack{\mathcal{P}\!\mathit{ic}}
\def\Picardfunctor{\mathrm{Pic}}
\def\Deformationcategory{\mathcal{D}\!\mathit{ef}}

% --- appended 2026-10-12 (Chief): the project own definitions of the reference/slogan/history
% display environments, copied mechanically from the shared Stacks preamble lines 8-12
% (\usepackage{verbatim} plus the three \newenvironment lines, which the chapter preamble keeps
% verbatim from the project preamble). Nothing added, removed or redefined: these environments
% carry no mathematical content and the project itself typesets them as comments. Without these
% three lines a counted statement whose own text uses them fails with Environment slogan undefined.
\usepackage{verbatim}
\newenvironment{reference}{\comment}{\endcomment}
\newenvironment{slogan}{\comment}{\endcomment}
\newenvironment{history}{\comment}{\endcomment}
\usepackage{cleveref}
\crefname{originallemma}{original lemma}{original lemmas}
\Crefname{originallemma}{Original lemma}{Original lemmas}
\makeatletter
\newlabel{lemma-c-square}{{0F16}{source}{}{sourcepointer-1}{}}
\newlabel{lemma-c-square@cref}{{[originallemma][0][]0F16}{[1][0][]source}{}{}{}}
\makeatother
\crefname{originallemma}{original lemma}{original lemmas}
\Crefname{originallemma}{Original lemma}{Original lemmas}
\makeatletter
\newlabel{lemma-square-triangle}{{0EVY}{source}{}{sourcepointer-2}{}}
\newlabel{lemma-square-triangle@cref}{{[originallemma][0][]0EVY}{[1][0][]source}{}{}{}}
\makeatother
\crefname{originallemma}{original lemma}{original lemmas}
\Crefname{originallemma}{Original lemma}{Original lemmas}
\makeatletter
\newlabel{lemma-K-injective-flat}{{093Y}{source}{}{sourcepointer-3}{}}
\newlabel{lemma-K-injective-flat@cref}{{[originallemma][0][]093Y}{[1][0][]source}{}{}{}}
\makeatother
\title{The blow-up square condition holds for complexes in the essential image of the derived pushforward}
\author{The Stacks Project Authors}
\date{stacks-project:sites-cohomology.tex:local-snapshot-20261011; GNU FDL 1.2 or later}
\begin{document}
\maketitle
\noindent\textit{Editorial note.} All delivered mathematical text is the author's own, copied from the pinned Stacks Project chapter file sites-cohomology.tex; no statement or proof step was rewritten, reordered or removed.\par\medskip
\noindent\textit{Editorial note.} The proof is the author's complete argument: a K-injective representative is chosen, the blow-up square lemma supplies the exact sequence of complexes, and the identification of that sequence with the one defining the canonical map is stated. The author's sentence 'we strongly urge the reader skip the proof' (source line 6781) is a reading suggestion concerning the brevity of the argument, not a statement that the proof is omitted, and it is preserved verbatim inside the counted span as required.\par\medskip
\noindent\textit{Editorial note.} Every labelled reference occurring inside the counted spans is carried in assembly.external\_references as a non-load-bearing navigation pointer (same-chapter ones at their definition site, cross-chapter ones at the invocation inside the counted span), so no load-bearing mathematics is replaced by a stub. No bibliography citation occurs inside the counted spans.\par\medskip
\noindent\textit{Editorial note.} The checked chapter list for same-name collisions is the set of stacks chapters staged in this repository: chow, cohomology, divisors, etale-cohomology, flat, more-algebra, more-morphisms, morphisms, spaces-more-morphisms, varieties together with this card's chapter sites-cohomology.\par\medskip
\noindent\textit{Editorial note.} The preamble is the chapter's own retained minimal preamble (sources/stacks-sites-cohomology/retained-preamble.tex), a verbatim excerpt of the native preamble: the Xy-pic setup and the native theorem declarations with the mathematical macros.\par\medskip
\noindent\textit{Allowed standard prerequisites (editorial).}\begin{itemize}
\item The morphism of sites epsilon and its basic properties. Same chapter, Section section-compare. Supplies the morphism whose derived pushforward is used.
\item The blow-up square produces a distinguished triangle. Same chapter, Lemma lemma-square-triangle. Supplies both the exactness of the displayed sequence of complexes and, through its proof, the identification of the resulting canonical map.
\item K-injectivity is preserved by the pushforward along epsilon. Same chapter, Lemma lemma-K-injective-flat. Shows that the pushforward of a K-injective representative represents the derived pushforward.
\end{itemize}\medskip
\ifdefined\hypertarget\hypertarget{sourcepointer-1}{}\fi
\noindent\textit{Source reference.} Original lemma 0F16: sources/stacks-sites-cohomology/pinned-original-source.tex, source lines 6001--6001. Navigation pointer invoked at proof line 6688; the counted span contains the pointer only and the cited result is not reproduced inside it, so the pointer is registered as a non-load-bearing navigation stub in the established form for Stacks material. No bibliography citation is involved.\par\medskip
\ifdefined\hypertarget\hypertarget{sourcepointer-2}{}\fi
\noindent\textit{Source reference.} Original lemma 0EVY: sources/stacks-sites-cohomology/pinned-original-source.tex, source lines 6074--6074. Navigation pointer invoked at proof line 6167; the counted span contains the pointer only and the cited result is not reproduced inside it, so the pointer is registered as a non-load-bearing navigation stub in the established form for Stacks material. No bibliography citation is involved.\par\medskip
\ifdefined\hypertarget\hypertarget{sourcepointer-3}{}\fi
\noindent\textit{Source reference.} Original lemma 093Y: sources/stacks-sites-cohomology/pinned-original-source.tex, source lines 4593--4593. Navigation pointer invoked at proof line 6786; the counted span contains the pointer only and the cited result is not reproduced inside it, so the pointer is registered as a non-load-bearing navigation stub in the established form for Stacks material. No bibliography citation is involved.\par\medskip
\noindent
Let $\mathcal{C}$ be a category. Let
$\text{Cov}(\mathcal{C}) \supset \text{Cov}'(\mathcal{C})$
be two ways to endow $\mathcal{C}$ with the structure of a site.
Denote $\tau$ the topology corresponding to $\text{Cov}(\mathcal{C})$
and $\tau'$ the topology corresponding to $\text{Cov}'(\mathcal{C})$.
Then the identity functor on $\mathcal{C}$ defines a morphism
of sites
$$
\epsilon : \mathcal{C}_\tau \longrightarrow \mathcal{C}_{\tau'}
$$
where $\epsilon_*$ is the identity functor on underlying presheaves and
where $\epsilon^{-1}$ is the $\tau$-sheafification of a $\tau'$-sheaf.

\begin{lemma}
\label{lemma-descent-squares-helper}
With $\epsilon : (\mathcal{C}_\tau, \mathcal{O}_\tau) \to
(\mathcal{C}_{\tau'}, \mathcal{O}_{\tau'})$ as above. Let
$$
\xymatrix{
E \ar[d] \ar[r] & Y \ar[d] \\
Z \ar[r] & X
}
$$
be a commutative diagram in the category $\mathcal{C}$ such that
\begin{enumerate}
\item $h_X^\# = h_Y^\# \amalg_{h_E^\#} h_Z^\#$, and
\item $h_E^\# \to h_Y^\#$ is injective
\end{enumerate}
where ${}^\#$ denotes $\tau$-sheafification. Then for
$K' \in D(\mathcal{O}_{\tau'})$ in the essential image of
$R\epsilon_*$ the map $c^{K'}_{X, Z, Y, E}$ of Lemma \ref{lemma-c-square}
(using the $\tau'$-topology) is an isomorphism.
\end{lemma}

\begin{proof}
This helper lemma is an almost immediate consequence of
Lemma \ref{lemma-square-triangle}
and we strongly urge the reader skip the proof.
Say $K' = R\epsilon_*K$. Choose a K-injective complex of
$\mathcal{O}_\tau$-modules $\mathcal{J}^\bullet$ representing $K$.
Then $\epsilon_*\mathcal{J}^\bullet$ is a K-injective complex of
$\mathcal{O}_{\tau'}$-modules representing $K'$, see
Lemma \ref{lemma-K-injective-flat}. Next,
$$
0 \to
\mathcal{J}^\bullet(X) \xrightarrow{\alpha}
\mathcal{J}^\bullet(Z) \oplus
\mathcal{J}^\bullet(Y) \xrightarrow{\beta}
\mathcal{J}^\bullet(E)
\to 0
$$
is a short exact sequence of complexes of abelian groups, see
Lemma \ref{lemma-square-triangle} and its proof.
Since this is the same as the sequence of complexes of abelian
groups which is used to define $c^{K'}_{X, Z, Y, E}$, we conclude.
\end{proof}

\end{document}
